% 386_Alpha_Geometrie_En_ch.tex
% Johann Pascher, 1 Oct 2026
% Where alpha hides: charge unit, two spheres and the reference energy

\providecommand{\BM}{\textbf{[B]}}
\providecommand{\KM}{\textbf{[K]}}
\providecommand{\SM}{\textbf{[S]}}
\providecommand{\XM}{\textbf{[X]}}
\providecommand{\QM}{\textbf{[Q]}}

\chapter*{Where \texorpdfstring{$\alpha$}{alpha} hides}
\addcontentsline{toc}{chapter}{Where alpha hides}

\begin{abstract}
FFGFT is first and foremost a ratio-based theory. It works in natural units with the convention
$\alpha=1$ and a correspondingly redefined charge; in this basic form it needs no fractal
correction. This document shows where $\alpha$ goes under the redefinition and what geometric
meaning it has there. With $\hbar=c=\varepsilon_0=1$ one has $e^2=4\pi\alpha$; with $\alpha=1$,
$e=\sqrt{4\pi}$, and the charge is measured by the full surface of the sphere. $\alpha$ is then
the square of the ratio of the historical to the geometric charge unit, and at the same time the
ratio of two spheres around the electron, $\alpha=r_e/\lambda_C$. The FFGFT bridge
$\alpha=\xi E_0^2$ reads $\xi E_0^2=1$ in T0 units, i.e. $E_0^2=1/\xi=7500$. The only opaque
element is the reference energy of the SI form $\alpha_{\text{SI}}=\xi(E_0/1\,\text{MeV})^2$: the
agreement says that this reference energy is one MeV to within $10^{-4}$. The MeV is not a geometric
quantity but a unit of the SI chain, built from the historical charge unit in which $\alpha$ also
hides.
For the fractal correction it follows that in the
$\alpha$ sector $K_{\text{frak}}$ appears only in the translation into SI; the factor required by
the two routes to $\alpha$ is the square of the bare reference energy in MeV.
\end{abstract}

% ============================================================
\section*{Status markers}
% ============================================================
\begin{center}
\begin{tabular}{@{}lp{11cm}@{}}
\textbf{[B]} & Algebraically proven. \\[2pt]
\textbf{[K]} & Numerically confirmed (check script, comparison with measured values). \\[2pt]
\textbf{[S]} & Assumption, open or speculative. \\[2pt]
\textbf{[X]} & Refuted or not sustainable. \\[2pt]
\textbf{[Q]} & Convention. \\
\end{tabular}
\end{center}
Constants: CODATA 2022. Check script: \texttt{pruef\_386\_alpha\_geometrie.py}.

% ============================================================
\section{The basic form of FFGFT}
% ============================================================

FFGFT calculates with ratios. In natural units, mass, energy, temperature and inverse time are
one quantity in four guises (A085), and the coupling is adopted as a unit of calculation:
$\alpha=1$ \QM\ (A267). The only parameter is $\xi=4/30000$. In this basic form there is no
fractal correction; $K_{\text{frak}}$, the SI values of $\alpha$ and $G$ and the fully worked-out
mass tables belong to the translation into SI, not to the core (A130). The question of this
document is therefore not how precisely $1/137$ is reproduced, but where $\alpha$ goes when it is
set to~1.

% ============================================================
\section{Where \texorpdfstring{$\alpha$}{alpha} goes under the redefinition}
% ============================================================

The fine-structure constant is
\begin{equation}
  \alpha=\frac{e^2}{4\pi\varepsilon_0\hbar c}.
\end{equation}
In natural units with $\hbar=c=\varepsilon_0=1$ (Heaviside--Lorentz) what remains is
\begin{equation}
  e^2=4\pi\alpha \quad\BM .
\end{equation}
The $4\pi$ is the full surface of the sphere. By Gauss's law the charge is the flux through a
closed sphere; the flux is spread over the solid angle $4\pi$. Setting $\alpha=1$ gives
\begin{equation}
  e_{\text{geo}}=\sqrt{4\pi}=3.545 ,
\end{equation}
so the charge is measured geometrically: its square is the full solid angle. With the measured
$\alpha$ the historical charge is $e_{\text{hist}}=\sqrt{4\pi\alpha}=0.3028$. Hence
\begin{equation}
  \alpha=\left(\frac{e_{\text{hist}}}{e_{\text{geo}}}\right)^2,\qquad
  \frac{e_{\text{geo}}}{e_{\text{hist}}}=\frac1{\sqrt\alpha}=11.706\quad\BM .
\end{equation}
Under the redefinition $\alpha$ therefore does not disappear. It moves into the charge unit:
$\alpha$ is the square of the ratio of the historical to the geometric charge unit. In every
calculation it stays on record through $e^2=4\pi\alpha$, one factor $\sqrt\alpha$ per vertex
(A267).

The same $4\pi$ appears in the Higgs--vacuum route: the factor $64\pi^4=16\pi^3\cdot4\pi$ of the
2025 vacuum formula is the $4\pi$ from $4\pi\varepsilon_0$ (Doc.~385).

% ============================================================
\section{The geometric picture: two spheres around the electron}
% ============================================================

Two spheres can be drawn around the electron. On the first, the Coulomb energy equals the rest
energy; this is the classical electron radius
\begin{equation}
  r_e=\frac{e^2}{4\pi\varepsilon_0\,m_ec^2}=2.818\times10^{-15}\ \text{m}.
\end{equation}
On the second, the quantum energy $\hbar c/r$ equals the rest energy; this is the reduced Compton
wavelength
\begin{equation}
  \lambda_C=\frac{\hbar}{m_ec}=3.862\times10^{-13}\ \text{m}.
\end{equation}
Their ratio is exactly the fine-structure constant:
\begin{equation}
  \frac{r_e}{\lambda_C}=\alpha ,\qquad \frac{a_0}{\lambda_C}=\frac1\alpha \quad\BM .
\end{equation}
Equivalently, $\alpha$ is the ratio of Coulomb energy to quantum energy on the same sphere,
$\bigl(e^2/(4\pi\varepsilon_0r)\bigr)/(\hbar c/r)$, for every radius~$r$ \KM.

Geometrically, $\alpha=1$ means that the Coulomb sphere and the Compton sphere coincide, and with
them the Bohr radius; $r_e$, $\lambda_C$ and $a_0$ are then a single scale $1/m_e$ (A267). The SI
value $1/137$ only says that in the historical charge unit the Coulomb sphere is 137 times smaller
than the Compton sphere. This is the transparent geometric meaning of $\alpha$.

% ============================================================
\section{The FFGFT bridge in ratio form}
% ============================================================

FFGFT links $\alpha$ to $\xi$ through the energy scale $E_0=\sqrt{m_em_\mu}$, the geometric mean
of the two lightest charged leptons. In the basic form, i.e. in T0 units with $\alpha=1$ and
without fractal correction, the bridge reads
\begin{equation}
  \xi\,E_0^2=1\quad\Longrightarrow\quad E_0^2=\frac1\xi=7500,\qquad E_0=86.60
  \quad\BM .
\end{equation}
The number $7500=1/\xi$ therefore belongs to $E_0^2$, not to $E_0$. Taken by itself, this
equation fixes the T0 energy unit: one unit is
$E_0/86.60=\sqrt{\xi\,m_em_\mu}=84.85$\,keV \KM. It acquires content only by comparison with a
second scale, and that is the transition to SI.

% ============================================================
\section{The SI bridge and the reference energy}
% ============================================================

In SI the ratio form becomes (Doc.~011, A130)
\begin{equation}
  \alpha_{\text{SI}}=\xi\left(\frac{E_0}{E_{\text{ref}}}\right)^2,\qquad E_{\text{ref}}=1\ \text{MeV}.
\end{equation}
This is where the bridge becomes opaque. $\xi E_0^2$ has the dimension energy$^2$, so a reference
energy is needed, and the corpus takes the MeV. Computing backwards which reference energy the
equation requires with the measured $\alpha$ gives
\begin{equation}
  E_{\text{ref}}=E_0\sqrt{\xi/\alpha}=
  \begin{cases}
  0.99323\ \text{MeV} & \text{bare, } E_0=7.348\ \text{MeV},\\
  0.99992\ \text{MeV} & \text{corrected, } E_0=7.397\ \text{MeV}
  \end{cases}\quad\KM .
\end{equation}
The bare value is $\sqrt{K_{\text{frak}}}$\,MeV to within $8\times10^{-5}$. The fractal
correction therefore enters only here, as a fine adjustment between the T0 unit and the MeV;
without it $\alpha^{-1}=7500/E_0^2=138.9$, with it $137.06$, $1.7\times10^{-4}$ from the measured
value. The same in the units of the basic form: the T0 energy unit is
$\sqrt\alpha\cdot E_{\text{ref}}$.

The agreement of $\alpha$ thus states one thing: the reference energy is one MeV to within
$10^{-4}$. This is a property of the anchor $E_0$ (R135), not an independent prediction \KM. What
the MeV is built from is shown in the next section.

% ============================================================
\section{Where the MeV comes from}
% ============================================================

The MeV is not a geometric quantity but a unit of the SI chain \QM. One has
\begin{equation}
  1\ \text{MeV}=10^6\cdot1.602176634\times10^{-19}\ \text{J}.
\end{equation}
The numerical factor is the numerical value of the elementary charge in coulombs. Since 2019 it has
been fixed; it was taken over from the old definition of the ampere via the force between two
conductors, with $\mu_0=4\pi\times10^{-7}$\,N/A$^2$. The joule depends on the kilogram, metre and
second. The MeV therefore carries the historical charge unit, exactly the unit in which $\alpha$
hides (Section~2).

The electron mass in MeV also has a measurement path in which $\alpha$ appears. What is measured
spectroscopically is the Rydberg energy $\mathrm{Ry}=13.6057$\,eV, and
\begin{equation}
  m_ec^2=\frac{2\,\mathrm{Ry}}{\alpha^2}=0.51100\ \text{MeV}\quad\BM .
\end{equation}
The SI form of the bridge, $\alpha=\xi\,m_em_\mu/(K_{\text{frak}}\,\text{MeV}^2)$, is equivalent to
\begin{equation}
  \frac{m_e}{1\ \text{MeV}}=\sqrt{\frac{\alpha\,K_{\text{frak}}}{\xi\,(m_\mu/m_e)}}=0.51104
  \quad(+8\times10^{-5})\quad\KM .
\end{equation}
In substance the relation thus fixes the ratio of the electron mass to the MeV through $\alpha$,
$\xi$ and the mass ratio. The MeV itself is a chain of fixed numbers.

The two spheres of Section~3 show on which side the correction belongs. The Compton sphere
$\lambda_C=\hbar/(m_ec)$ contains only the mass, i.e. the anchor. The Coulomb sphere
$r_e=e^2/(4\pi\varepsilon_0m_ec^2)$ contains $e$, i.e. the same historical unit from which the MeV
is built. If the fractal correction is placed on this side instead of on the masses, the Compton
sphere stays bare and the reference energy becomes
\begin{equation}
  E_{\text{ref}}=\frac{E_0^{\text{bare}}\sqrt{\xi/\alpha}}{\sqrt{K_{\text{frak}}}}
  =0.99992\ \text{MeV},
\end{equation}
i.e. $1$\,MeV to within $8\times10^{-5}$. Mathematically this is the same equation as with the
corrected $E_0$ (R72); the reading is more coherent, however, because the masses stay bare and
$K_{\text{frak}}$ sits on the side of the unit, i.e. in the SI translation. This assignment is a
choice \SM. The trail explains what the MeV is built from, not why it fits to within $10^{-4}$; the
residual of $8\times10^{-5}$ remains open \SM.

% ============================================================
\section{Where the fractal correction sits}
% ============================================================

The reference energy shows where $K_{\text{frak}}$ sits in the $\alpha$ sector. The two routes
to $\alpha$ in A130 differ only by this factor: route~1 inserts the bare mean
$E_0=\sqrt{m_em_\mu}$ and gives $\alpha^{-1}=138.9$, route~2 the corrected one and gives
$137.06$. The factor that $\alpha$ requires between the two is
\begin{equation}
  K_\alpha=\xi\,m_em_\mu\,\alpha^{-1}=\left(\frac{E_{\text{ref}}^{\text{bare}}}{1\ \text{MeV}}\right)^2
  =0.98650\quad\BM ,
\end{equation}
$1.7\times10^{-4}$ from $74/75$ \KM. This is the same equation as in A130, only read differently,
and therefore not a new measurement. It does show where the factor stands: $K_\alpha$ is the
square of the bare reference energy measured in MeV, i.e. exactly the ratio between the T0 energy
unit and the MeV. In the basic form it does not occur: there the bridge reads $\xi E_0^2=1$, which
fixes the unit without leaving room for a correction factor \BM.

This agrees with Doc.~122: in ratios $K_{\text{frak}}$ cancels, it acts only on absolute
quantities measured against an external reference. For the $\alpha$ sector it follows that
$K_{\text{frak}}$ appears only in the translation into SI, as the adjustment between the T0 unit
and the MeV \KM. Its value $74/75=1-100\xi$, by contrast, comes from the geometry, as the rotation
number of the frozen recursion (Docs.~381, 384) and from the cumulative running (A040); origin and
appearance are therefore separate. The residual of $1.7\times10^{-4}$ between $K_\alpha$ and
$74/75$ remains open (R135, R139) \SM.

% ============================================================
\section{What holds and what is open}
% ============================================================

\begin{center}
\small
\begin{tabular}{@{}p{7.4cm}p{1.2cm}p{5.2cm}@{}}
\toprule
\textbf{Statement} & \textbf{Status} & \textbf{Reason} \\
\midrule
$e^2=4\pi\alpha$; with $\alpha=1$, $e=\sqrt{4\pi}$ & \BM & $\hbar=c=\varepsilon_0=1$ \\
$\alpha=(e_{\text{hist}}/e_{\text{geo}})^2$ & \BM & redefinition of the charge \\
$\alpha=r_e/\lambda_C$, $a_0/\lambda_C=1/\alpha$ & \BM & two spheres around the electron \\
$\alpha=1$ as unit of calculation & \QM & A267 \\
$\xi E_0^2=1$, $E_0^2=1/\xi=7500$ in T0 units & \BM & basic form \\
T0 energy unit $\sqrt{\xi m_em_\mu}=84.85$\,keV & \KM & check script \\
$E_{\text{ref}}=0.99992$\,MeV (corrected), $0.99323$\,MeV (bare) & \KM & property of the anchor (R135) \\
$1$\,MeV $=10^6\cdot1.602\times10^{-19}$\,J, numerical value of $e$ & \QM & SI chain, $\mu_0=4\pi\times10^{-7}$ \\
$m_ec^2=2\,\mathrm{Ry}/\alpha^2$ & \BM & measurement path of the electron mass \\
$m_e/\text{MeV}=\sqrt{\alpha K_{\text{frak}}/(\xi\,m_\mu/m_e)}=0.51104$ & \KM & $+8\times10^{-5}$, same residual \\
correction on the unit side, Compton sphere bare & \SM & reading; residual $8\times10^{-5}$ open \\
$K_\alpha=(E_{\text{ref}}^{\text{bare}}/\text{MeV})^2=0.98650$ & \BM & same equation as A130 \\
$K_{\text{frak}}$ in the $\alpha$ sector only in the SI translation & \KM & basic form $\xi E_0^2=1$; Doc.~122 \\
residual $1.7\times10^{-4}$ between $K_\alpha$ and $74/75$ & \SM & R135, R139 \\
\bottomrule
\end{tabular}
\end{center}

% ============================================================
\section{Aligned places in the corpus}
% ============================================================

With this document the places with $E_0=1/\xi$ carry a dated note pointing to the ratio form
$E_0^2=1/\xi$: Doc.~005 (previously without a note), 032, 133 and 165, each De/En. Doc.~384 puts
the basic form first. Doc.~013, like Doc.~014 before it, notes that $S_{T0}=1\,\text{MeV}/c^2$ holds by construction. A130 notes that the gap between the two routes is the square of the bare
reference energy.

% ============================================================
\section{Conclusion}
% ============================================================

$\alpha$ does not disappear under the redefinition; it becomes the charge unit. With $\alpha=1$
the charge is measured by the full surface of the sphere, $e=\sqrt{4\pi}$, and the Coulomb sphere
of the electron coincides with its Compton sphere. The measured value $1/137$ is the radius ratio
$r_e/\lambda_C$ in the historical charge unit. In its basic form the FFGFT bridge reads
$\xi E_0^2=1$, i.e. $E_0^2=1/\xi$; this is transparent. Only the step to SI is opaque, because it
needs a reference energy that is one MeV to within $10^{-4}$. The MeV is built from the numerical
value of the historical charge, i.e. from the same unit in which $\alpha$ hides. On this side, and
only there, the fractal correction sits: the Compton sphere stays bare, the value $74/75$ comes from
the geometry, and it takes effect only in the adjustment to the MeV.

\vspace{1em}
\noindent\textit{Check script:} \texttt{2/python/Dok386\_Skripte/pruef\_386\_alpha\_geometrie.py}
--- 21/21.
